% BEGIN VOL08-EXPANSION VIII27
\section{Supplementary worked examples}
\label{sec:viii27-pedagogy}

\begin{example}[The torus from two circles]\label{ex:viii27-ped-01}
Since $H_*(S^1)$ is free, the Tor terms vanish in Kunneth. Thus
\[
H_n(S^1\times S^1)\cong\bigoplus_{p+q=n}H_p(S^1)\otimes H_q(S^1),
\]
giving $\mathbb Z$, $\mathbb Z^2$, $\mathbb Z$ in degrees $0,1,2$.
\end{example}

\begin{example}[A product of spheres]\label{ex:viii27-ped-02}
For $S^m\times S^n$ with $m,n>0$, free homology gives classes in degrees $0,m,n,m+n$ (with merging when $m=n$). The top class is the cross product of the two sphere fundamental classes.
\end{example}

\begin{example}[Where Tor enters]\label{ex:viii27-ped-03}
For spaces with torsion homology, Kunneth contains a shifted Tor term. This explains why product homology may contain classes not visible in the simple tensor products of same-total-degree homology groups.
\end{example}

\section{Graded supplementary exercises}
The following exercises progress from direct computation and construction to proof, hypothesis testing, application, and synthesis.

\subsection*{Standard computations and constructions}

\begin{exercise}[Product with a point]\label{exr:viii27-ped-01}
Compute $H_*(X\times\{*\})$ using Kunneth.
\end{exercise}
\begin{hint}
$H_0(\{*\})=\mathbb Z$ is the only nonzero factor.
\end{hint}
\begin{solution}
One recovers $H_*(X)$.
\end{solution}

\begin{exercise}[Torus H1]\label{exr:viii27-ped-02}
Compute $H_1(S^1\times S^1)$.
\end{exercise}
\begin{hint}
Use degrees $(1,0)$ and $(0,1)$.
\end{hint}
\begin{solution}
$H_1\cong\mathbb Z\oplus\mathbb Z$.
\end{solution}

\begin{exercise}[Torus H2]\label{exr:viii27-ped-03}
Compute $H_2(T^2)$.
\end{exercise}
\begin{hint}
Use the tensor $H_1(S^1)\otimes H_1(S^1)$.
\end{hint}
\begin{solution}
$H_2(T^2)\cong\mathbb Z$.
\end{solution}

\begin{exercise}[Sphere product top class]\label{exr:viii27-ped-04}
Compute $H_{m+n}(S^m\times S^n)$.
\end{exercise}
\begin{hint}
Only the product of top classes contributes.
\end{hint}
\begin{solution}
It is $\mathbb Z$.
\end{solution}

\begin{exercise}[Euler product]\label{exr:viii27-ped-05}
Use Kunneth ranks to recover $\chi(X\times Y)=\chi(X)\chi(Y)$ when homology is finite and free.
\end{exercise}
\begin{hint}
Alternate the tensor-product ranks.
\end{hint}
\begin{solution}
The double alternating sum factors as the product of the two Euler-characteristic sums.
\end{solution}

\subsection*{Proofs}

\begin{exercise}[Free-factor Kunneth]\label{exr:viii27-ped-06}
Prove Tor vanishes if one factor has free integral homology.
\end{exercise}
\begin{hint}
Tor with a free abelian group is zero.
\end{hint}
\begin{solution}
Every homology group of that factor is free, so all Tor summands in the Kunneth sequence vanish.
\end{solution}

\begin{exercise}[Cross-product generators]\label{exr:viii27-ped-07}
Explain how classes $a\in H_p(X)$ and $b\in H_q(Y)$ produce a class in $H_{p+q}(X\times Y)$.
\end{exercise}
\begin{hint}
Use the homology cross product.
\end{hint}
\begin{solution}
The Eilenberg--Zilber map sends the tensor of representative cycles to a cycle on the product, defining $a\times b$.
\end{solution}

\begin{exercise}[Naturality]\label{exr:viii27-ped-08}
State the naturality of cross products under maps $f:X\to X'$ and $g:Y\to Y'$.
\end{exercise}
\begin{hint}
Apply the product map.
\end{hint}
\begin{solution}
$(f\times g)_*(a\times b)=f_*(a)\times g_*(b)$.
\end{solution}

\begin{exercise}[Kunneth short exact structure]\label{exr:viii27-ped-09}
Explain why the Tor term is shifted by one total degree.
\end{exercise}
\begin{hint}
Recall the standard homological Kunneth exact sequence.
\end{hint}
\begin{solution}
The Tor contribution to $H_n(X\times Y)$ comes from pairs with $p+q=n-1$, reflecting the first derived functor of tensor product.
\end{solution}

\subsection*{Counterexamples and hypothesis tests}

\begin{exercise}[Always tensor only]\label{exr:viii27-ped-10}
Test: $H_n(X\times Y)$ is always just $\bigoplus_{p+q=n}H_p(X)\otimes H_q(Y)$.
\end{exercise}
\begin{hint}
Torsion can create Tor terms.
\end{hint}
\begin{solution}
False over $\mathbb Z$ unless suitable freeness hypotheses remove Tor.
\end{solution}

\begin{exercise}[Product Betti numbers add]\label{exr:viii27-ped-11}
Test: $b_n(X\times Y)=b_n(X)+b_n(Y)$.
\end{exercise}
\begin{hint}
Over a field, use convolution.
\end{hint}
\begin{solution}
False. Betti numbers satisfy $b_n(X\times Y)=\sum_{p+q=n}b_p(X)b_q(Y)$.
\end{solution}

\begin{exercise}[Top class without orientability]\label{exr:viii27-ped-12}
Test: every product of closed manifolds has an integral top class.
\end{exercise}
\begin{hint}
Orientability matters.
\end{hint}
\begin{solution}
False; an integral fundamental class requires orientability of the product, equivalently of both factors in the connected case.
\end{solution}

\subsection*{Applications and investigations}

\begin{exercise}[High-dimensional construction]\label{exr:viii27-ped-13}
How can Kunneth compute a large product without constructing its boundary matrix?
\end{exercise}
\begin{hint}
Use already-known homology of the factors.
\end{hint}
\begin{solution}
It replaces a potentially huge product chain complex by tensor and Tor operations on much smaller homology groups.
\end{solution}

\begin{exercise}[Persistent feature intuition]\label{exr:viii27-ped-14}
Why is the cross product useful when interpreting independent topological features in two factors?
\end{exercise}
\begin{hint}
A feature in each factor combines into a higher-dimensional product feature.
\end{hint}
\begin{solution}
A $p$-cycle and a $q$-cycle can form a $(p+q)$-cycle, such as two loop classes combining into the torus fundamental class.
\end{solution}

\subsection*{Challenge problems}

\begin{exercise}[Three-torus]\label{exr:viii27-ped-15}
Compute the integral Betti numbers of $T^3=(S^1)^3$.
\end{exercise}
\begin{hint}
Iterate Kunneth or expand $(1+t)^3$.
\end{hint}
\begin{solution}
The Betti numbers are $1,3,3,1$.
\end{solution}

\begin{exercise}[Kunneth synthesis]\label{exr:viii27-ped-16}
Describe a practical Kunneth workflow over $\mathbb Z$.
\end{exercise}
\begin{hint}
List tensor contributions, then Tor contributions, then resolve the short exact sequence.
\end{hint}
\begin{solution}
For each total degree, compute all $H_p(X)\otimes H_q(Y)$ with $p+q=n$, all Tor groups with $p+q=n-1$, assemble the Kunneth short exact sequence, and use any available splitting or structural information.
\end{solution}

% END VOL08-EXPANSION VIII27
